% 本文件由 LaTeX 排版 Agent 根据有效 translation.json 自动生成。
% author=John Chrysostom; work_id=1904; title=致一位年轻寡妇的书信

\chapter{致一位年轻寡妇的书信}

% structure: p0001 source=h1 target=omitted reason=page-title-already-rendered-by-parent

以下书信的日期可确定在很窄的范围内。信中提及（第五章）\Person{瓦伦斯}皇帝在公元378年与哥特人于\Place{阿德里安堡}的战役中败亡，作为近期事件。而被描述为自即位以来不断从事战争的皇帝（第四章）必是继瓦伦斯之后的\Person{狄奥多西}；由于哥特人仍肆意蹂躏大片区域，并傲慢地嘲笑帝国军队的胆怯（同上），此信必写于狄奥多西在382年给予他们沉重打击之前。整封信深深浸透着对人生不幸与无常的深切感受，这是社会的道德败坏和帝国近期的灾难以特殊力量刻印在人们心中的，常在异教徒中产生犬儒式的阴郁或漠不关心，却引导基督徒更紧切地坚守福音的盼望与安慰。\par

1. 你受到了沉重的打击，且从上而来的武器刺入了要害，这一点众人都愿承认，即使最严苛的伦理学家也不会否认；然而，被忧伤所伤的人不应将全部时间用于哀悼和眼泪，而应为医治伤口做好充分准备，免得因忽视而让眼泪加剧创伤，让忧伤之火愈燃愈烈；故此，倾听安慰之言是好事，至少暂时抑制你的泪泉，将自己交托给那些努力安慰你的人。为此，当你的忧伤达到顶峰、雷霆刚刚降临时，我忍住不打扰你；而是等候一段时间，让你尽情哀哭；如今你已能稍从迷雾中观望，并向那些试图安慰你的人敞开耳朵，我也愿藉自己的贡献来附和你的婢女们的话语。因为当风暴仍然猛烈、忧伤的狂风劲吹时，劝人停止悲伤只会激起更甚的哀嚎，招致仇恨，反而给火焰添柴，且被视为冷酷愚昧之人。但当汹涌的水面开始平息，天主缓和了波涛的狂怒时，我们便可自由地展开言谈之帆。因为在适度的风暴中，技艺或许能发挥作用；但当风势不可抗拒时，经验便无济于事了。因此，我迄今保持沉默，直到现在才敢打破沉默，因为从你的叔父处听说，可以开始鼓起勇气，因你一些更可敬的婢女如今敢于详细谈论这些事，此外还有你家以外的妇女，或是你的亲属，或是以其他资格胜任此职。现在，你若允许她们与你交谈，我怀有极大的希望和信心，你必不会轻视我的言语，而会尽量平静安宁地倾听。无论如何，女性性别的确更易感受痛苦；若再加上年轻、过早的寡居、缺乏处事经验、众多忧虑的纷扰，以及先前整个生活在奢华、欢乐和财富中被滋养，邪恶便成倍增加；若她遭受此苦而得不到从上而来的帮助，即便偶然的思绪也能使她崩溃。我认为这是天主眷顾你首要且最大的证据：因为当如此巨大的患难突然同时袭来折磨你时，你没有被忧伤压倒，也没有失去本性的心态，这并非出自任何人的帮助，而是出自全能的手——其智慧不可测度，其明智无法探究，即\BibleQuote{仁慈之父和一切安慰的天主}\BibleRef{格后 1:3}。因为经上说：\BibleQuote{原来祂打伤了我们，也必包扎；祂击打我们，也必治疗，使我们痊愈。}\BibleRef{欧 6:2}\par

因为当你那蒙福的丈夫还在你身边时，你享有尊荣、关怀和热切的照顾；实际上，你享有你从丈夫那里所期待的一切；但自从天主将他接到自己身边后，祂就接替了他对你的位置。这不是我的话，而是蒙福先知\Person{达味}的话，因为他说：\BibleQuote{祂必抚养孤儿和寡妇}，在另一处又称祂为\BibleQuote{孤儿的慈父、寡妇的判官}；因此，在许多经文中你将看到祂认真关注这一类人的处境。\par

2. 但是，为了避免\Quote{寡妇}这个称呼的不断重复搅扰你的灵魂，扰乱你的理智（这称呼在你青春正盛时加诸于你），我先要就此点加以论述，并向你证明：\Quote{寡妇}这个称呼不是灾祸的称号，而是荣耀的称号，甚至是最高的荣耀。因为请不要引用世俗的错误意见作为证据，而要引用真福保禄的劝诫，或更确切地说，基督的劝诫。因为在他的话语中，基督藉着他说话，正如他自己所说：\BibleQuote{若你们寻求在我内说话的基督的凭据吗？}\BibleRef{格后 13:3}那么他说什么呢？\BibleQuote{不要录用一个六十岁以下的寡妇}\BibleRef{弟前 5:9}又在另处说：\BibleQuote{至于年轻的寡妇，你要拒绝}\BibleRef{弟前 5:11}藉着这两句话，他向我们表明这事的重要性。当他为主教制定规则时，他并未规定年龄标准，但在此事上他却非常讲究，请问这是为什么呢？不是因为寡居高于司祭职，而是因为寡妇比司铎要承担更大的劳苦，被许多公私事务四面围攻。因为一座没有城墙的城市暴露在一切想掠夺它的人面前，同样，一个年轻的寡妇生活在寡居中，四面八方都有许多人图谋她，不仅那些想获取她钱财的人，还有那些一心要败坏她贞洁的人。除此之外，我们还会发现她还要承受其他可能导致她跌倒的情况。因为仆人的轻视、对事务的疏忽、昔日所受尊敬的丧失、看见同龄人兴旺发达，以及对奢侈生活的渴望，都诱使妇女再婚。有些人并不选择通过婚姻法律与男人结合，而是暗中秘密行事。他们这样做是为了享有寡居的赞誉；可见这状态似乎并不受人责备，反而受人钦佩，并被视为在人眼中值得尊敬的，不仅在我们信徒中，即使在外教人中也是如此。因为有一次，当我还是年轻人时，我记得教导我的那位诡辩家（他在敬神方面超越众人）在大庭广众之下对我的母亲表示钦佩。他照常向坐在他旁边的人询问我是谁，有人告诉他说，我是一个寡妇的儿子，他就问我母亲的年龄和寡居多久了，我告诉他她四十岁，其中二十年是在失去我父亲后度过的，他惊讶地大声喊道，转向在场的人说：\Quote{天哪！基督徒中竟有这样的女子！}寡居所享有的赞美和钦佩，不仅在我们中间，也在教会之外的人中，竟是如此之大。圣保禄深知这一切，所以他说：\BibleQuote{不要录用一个六十岁以下的寡妇}。即使过了这个年龄的严格要求，他也不允许她列入这一神圣团体，而是提出一些附加条件：\BibleQuote{因善行而获有声望，如教养儿女，款待旅客，给圣徒洗脚，救济遭难的人，勤行各种善工}\BibleRef{弟前 5:10}。天哪！这是怎样的考验和审视！他对寡妇要求多少德行，又是何等精确地界定了它！若他无意将一种荣耀与尊贵的地位托付给她，他就不会这样做了。\BibleQuote{至于年轻的寡妇}他说\BibleQuote{你要拒绝}，然后他加上理由：\BibleQuote{因为当他们放纵情欲，违背基督的时候，便想再嫁}\BibleRef{弟前 5:11}。藉着这话，他使我们明白，那些失去丈夫的人，基督代替了她们的丈夫。请注意他如何藉着表明这结合的温和与轻松性质来断言这一点；我指的是那处经文：\BibleQuote{因为当他们放纵情欲，违背基督的时候，便想再嫁}，好像他是一位温和的丈夫，不行使权柄管辖她们，而是让她们自由生活。保禄不仅在此处谈论了此事，而且在另一处地方也表明了对此的深切关怀，他说：\BibleQuote{至于那纵情享乐的寡妇，虽生犹死；而那真正寡居而孤独的，寄望于天主，日夜在恳求和祈祷中坚持}\BibleRef{弟前 5:6}。他给格林多人写信时也说：\BibleQuote{可是，如果她这样守下去，她更为有福}\BibleRef{格前 7:40}。你看，寡居受到了多么大的赞美，这还是在《新约》中，当时贞洁的美德也已清楚显现。然而，即使这状态的光彩也不能遮蔽寡居的光荣，它依然明亮地闪耀，保持其自身的价值。因此，当我们不时提到寡居时，不要沮丧，也不要认为此事是耻辱；如果这是耻辱，那么贞洁更是耻辱。但事实并非如此；不！天主不许。既然我们所有人都钦佩并欢迎那些在丈夫在世时守节的妇女，难道我们不应同样为那些在丈夫去世后依然对丈夫怀有同样美好感情并相应称赞她们的妇女而感到喜悦吗？正如我那时所说，当你与可敬的特拉修斯同住时，你享有妻子从丈夫那里自然得到的尊荣和眷顾；但如今，代替他的是万有的主天主，祂从前就是你的保护者，如今更将更加热切地保护你；正如我已经说过的，祂已显示了一个不小的庇护凭据，即在如此焦虑和悲伤的炉火中保全你完好无损，没有让你遭受任何不愉快的事。既然祂不曾容许在如此汹涌的水中发生任何海难，更何况在风平浪静时，祂必定会保全你的灵魂，减轻你寡居的重担以及那些看似可怕的后果。\par

3. 但若使你痛苦的并非寡妇的名分，而是失去如此一位丈夫——我承认，全世界从事世俗事务的人中，少有像他那样深情、温柔、谦卑、真诚、善解人意、虔诚的人。诚然，若他完全消亡、不复存在，那么忧伤悲恸自是应当；但若他只是驶入了宁静的港湾，启程前往那真正是他君王的天主那里，就不应哀悼，反该为此欢喜。因为这样的死亡并非死亡，而只是一种迁徙和迁移，从较坏到较好，从地上到天上，从人间到天使、总领天使，以及天使和总领天使的主那里。因为在此世，当他服务皇帝时，常预期有危险，且有许多心怀恶意之人所策划的阴谋，因为他的声望越高，敌人的计谋也越多；但现在他离世到了另一世界，这些事都无须担忧了。因此，你既为天主带走如此良善可敬的人而悲伤，也该为他平安尊荣地离世，且从困扰这危险时节的烦恼中解脱，居于极大的平安与宁静而欢喜。因为承认天堂远胜于地上，却仍为那些从此世转移到彼世的人哀悼，岂非不合情理？因为如果你那蒙福的丈夫曾是那些生活无耻、违背天主所喜悦之事的人之一，那么不仅在他离世后，即使他还活着时，也当为他哀哭叹气；但既然他是天主的朋友之一，我们不仅应在他在世时，也应在他人息安眠后为他欢喜。我们理应如此行，你必已听过蒙福的保禄的话：\Quote{离去并与基督同处，远为更好。}\BibleRef{斐 1:33}但也许你渴望听到你丈夫的话语，享受你曾倾注于他的情感，你思念他的陪伴，以及因他而有的荣耀、辉煌、尊荣、安稳；这一切都已消逝，令你忧伤并黯淡了你的生命。嗯！你曾倾注于他的情感，如今仍可像从前一样保持。\par

因为爱的力量就是如此：它不仅拥抱、联合、紧固那些在眼前、临近、可见的人，也拥抱那些遥远的人；无论是时间的长久、空间的分离，还是其他任何事，都不能打破和撕裂灵魂的情感。但若你渴望与他面对面相见（我知道这是你特别向往的），就为他的尊荣保持你的床榻神圣，不受任何其他男人的触碰，并竭力活出像他一样的生命，那么有一天你必定会离去，加入与他的同一群体，不是像在此世那样与他同住五年，也不是二十年、一百年，也不是一千年或两千年，而是无穷无尽的世代。因为使人有资格继承那些安息之地的，并非任何肉身关系，而是生活方式的一致。因为如果是道德品格的同一性使拉匝禄——虽与亚巴郎非亲非故——与他同在天上的怀抱中，并使许多从东从西而来的人得以与他同席，那么安息之地也必将接纳你和那位善\Person{德辣西乌斯}，只要你表现出与他相同的生活方式，那时你将重新得回他，不再是他在世时那样的肉身之美，而是另一种光彩，太阳光芒也为之逊色的辉煌。因为此肉身，即使达到极高的美貌标准，仍然是可朽的；但那些中悦天主之人的身体，将被赋予一种连这些眼睛都无法直视的荣耀。天主给了我们某些标记和模糊的征兆，无论是在旧约还是在新约时代。因为在前者中，梅瑟的面容闪耀着令以色列子民眼睛无法承受的光辉；而在新约中，基督的面容比梅瑟的更加闪耀。请告诉我，若有人应许让你丈夫作全地的王，然后命你为他退隐二十年，并许诺之后将头戴冠冕、身着紫袍的他归还给你，使你再次与他同列，你难道不会以适当的自制温顺地忍受这分离吗？你难道不会为此恩赐而欣喜，并视其为值得祈求的事吗？那么，现在就顺从吧，不是为了地上的国度，而是为了天上的国度；不是为了得回他穿着金色衣袍，而是为了他披上不朽与荣耀，正如天上的居民所应得的那样。若你因考验历时长久而觉得难以承受，或许他会在异象中探望你，像往常一样与你交谈，并让你看到你渴望的容颜：愿这成为你的慰藉，代替书信，尽管事实上它比书信更为真切。因为后者不过是些笔迹线条供你观看，而前者你却能看到他的面容轮廓、温柔微笑、身形与动作，听见他的言谈，认出你所深爱的那声音。\par

4. 但既然你也为失去从前因他而享有的安稳而悲伤，或许还为了那些即将实现的显赫希望（因为我曾听说他很快就要升任总督，我想这正是尤其扰乱和折磨你灵魂的事），我恳求你想想那些官位比他更高，却落得非常悲惨结局的人。让我提醒你：你大概听说过\Person{西西里的狄奥多禄}其人吧，他是最杰出的人物之一，在体魄、美貌以及皇帝对他的信任方面都超过众人，而且比王室任何成员都更有权势，但他并未谦卑地承受这好运，因参与一项反对皇帝的阴谋而被俘，惨遭斩首；他的妻子在教育、出身和其他各方面丝毫不逊于你这高贵之人，却突然被剥夺了一切所有，甚至被剥夺了自由，被编入家奴之列，被迫过着比任何婢女都更可怜的生活，只有这一点胜过其余：由于她灾难的极端严重，使所有看见她的人都流下眼泪。据说\Person{阿尔忒弥西亚}也是一位声名显赫之人的妻子，因她丈夫也图谋篡位，而陷入同样的贫困境地，甚至失明；因为绝望之深和泪水之多毁坏了她的视力；如今她需要人搀扶她的手，领她到别人门口去获取必要的食物。我还可以提到许多其他这样败落的家族，但我知道你天性虔诚明智，不愿从别人的不幸中为自己的灾难寻求安慰。我刚才提到那些例子的唯一原因，是让你明白：人的事物都是虚无，而先知的话千真万确：\BibleQuote{凡人的光荣尽如草上的花}。\BibleRef{依 40:5} 因为人越是尊崇高贵，他们所遭受的毁灭也越大，这不仅适用于被统治的人，也适用于统治者本身。因为不可能找到任何一个私人家庭曾陷入如此巨大的灾难，就像皇室所浸染的祸患那样。因为父母和配偶的早逝，以及比悲剧中更残忍、更痛苦的暴力死亡，特别困扰着这种政体。\par

现在略过古代，在我们这一代执政的人中，总共九位，只有两人得以寿终正寝；其余人中，一个被篡位者所杀，一个死于战场，一个被卫队阴谋杀害，一个被选举并授他紫袍的人亲手所杀；他们的妻子，有的据报死于毒药，有的因悲伤过度而亡；而那些仍然活着的妻子中，一位带着孤儿儿子，惊恐不安，唯恐当权者中有人因惧怕将来可能发生的事而将他除掉；另一位勉强应允多番恳求，从被掌权者驱逐的放逐中返回。至于现今统治者的妻子，那位从先前灾难中稍得恢复的，欢乐中混杂着许多悲伤，因为权力拥有者尚年轻无经验，周围有许多图谋不轨之人；另一位则恐惧欲死，比被判死刑的罪犯还要痛苦，因为她的丈夫自加冕以来直到今日，始终不断征战，四面受辱和斥责，比实际的灾难更使他憔悴。因为从未有过的事如今发生了：蛮族离开本国，无数次侵占了我国无限的土地，放火烧地，攻占城镇，却不打算返回家园，反而以过节而非打仗的姿态，对我们所有人嗤之以鼻；据说他们的一个国王宣称，他对我方士兵的厚颜无耻感到惊讶，他们被杀比羊还容易，却仍期望得胜，不愿离开自己的国家；他说他自己已经厌倦了砍杀他们的工作。试想皇帝及其皇后听到这些话时是何感受！\par

5. 既然我提到了这场战争，我便想起了一大群寡妇，她们从前因丈夫的尊荣而大放异彩，如今却全都穿上深色丧服，终日哀叹。因为她们并未享有你亲爱的自己所享的福分。你这卓越的朋友，你曾看见你那良善的丈夫躺在病床上，听见他的临终遗言，接受他关于家庭事务的指示，并得知通过他的遗嘱如何防范贪婪和诡诈之人的各种侵犯。不仅如此，当他已躺卧去世时，你多次扑在他身上，亲吻他的眼睛，拥抱他，为他哀哭，你看见他以极大尊荣被送往埋葬，按礼节做了一切必要的殡葬之事，并因常去他的坟墓而获得了不少安慰。但这些妇女却被剥夺了一切：她们满怀希望将丈夫送上前线，盼着他们归来，却只接到他们战死的噩耗。没有人将阵亡者的遗体带回给她们，也没有带回任何东西，只带来他们死状的描述。有些人甚至没有得到这样的消息，也不知道丈夫是如何倒下的，因为他们被埋在战场上的阵亡者堆中。\par

若将军们多半如此丧命，又有什么可奇怪的呢？连皇帝本人带着少数士兵被困在一个村庄里，不敢出去迎敌，反而留在屋内；当敌人纵火烧房时，他便与里面所有的人一同烧死——不单是人，还有马匹、梁木和墙壁，整个化为灰烬。这就是那些随皇帝出征的人归来时，代替皇帝本人带回给他妻子的事迹。因为世界的荣华，与舞台上所发生的事以及春天的花朵之美，毫无区别。首先，它们尚未显现便已消逝；其次，即使能持续片刻，也很快衰败。众人的尊敬和荣耀，有什么比这更无价值的呢？它有什么果实？有什么益处？能满足什么有用的目的？但愿这仅仅是祸患而已！但实际上，除了从拥有中得不到任何好处外，拥有这最残酷的女主人的人，还不断被迫承受许多痛苦和有害之事；因为她是那些拥有她的人的女主人，她越是被她的奴隶奉承，就越是高抬自己，用越加严厉的命令捆绑他们；但她永远无法向那些轻视和忽略她的人报复；她比任何暴君和野兽都更凶猛。因为暴君和野兽常因讨好的态度而变得温和，但她的怒火却在我们对她最逢迎时最为强烈；如果她发现有人听从她、事事顺从她，她就绝不会再停止发号施令。此外，她还有一个盟友，称之为她的女儿也不为过。因为当她自身成熟并深深扎根在我们中间之后，她便产生了傲慢，这东西与她自己一样，能将拥有它之人的灵魂推向彻底的毁灭。\par

6. 那么请告诉我，你之所以哀恸，是否因为天主将你从如此残酷的束缚中解救出来，并堵住了这些致命疾病的所有通道？因为当你丈夫在世时，它们不断攻击你心中的思想；但他去世后，它们便再无立足之处来抓住你的悟性。因此，未来的修练应当是——戒除为这些祸患的消失而哀恸，戒除渴望它们所施行的痛苦暴政。因为当它们猛烈吹刮时，便从根基推翻一切，将它们打得粉碎；正如许多妓女，虽然天生丑陋，却用彩釉和颜料挑动年轻人尚柔嫩的情感，一旦将他们控制，便比对待任何奴隶都更加蛮横；同样，这些情欲——虚荣和傲慢——比任何其他污染更玷污人的灵魂。\par

因此，财富在大多数人眼中也似乎是一件好事；至少当它被剥离了这种虚荣的激情时，它就不再显得可取了。无论如何，那些被允许在贫穷中获得大众荣耀的人，不再偏好财富，反而鄙视大量赐予他们的黄金。你不需要从我这里学习这些人是哪些，因为你比我更了解他们：\Person{埃帕米农达}、\Person{苏格拉底}、\Person{阿里斯提德}、\Person{第欧根尼}、\Person{克拉特斯}，后者把自己的土地变成了牧场。诚然，其他人，既然他们不可能致富，看到荣耀在贫穷中带给他们，就立刻投身于荣耀；但这个人甚至抛弃了他所拥有的东西；他们是如此痴迷于追逐这残酷的怪物。那么，我们不要哭泣，因为天主已将我们从这可耻的奴役中解救出来，这奴役是嘲笑和许多责备的对象；因为它本身没有什么辉煌，除了它拥有的名字，而实际上它使拥有它的人处于一种与其名称不符的境地，没有人不嘲笑任何为了荣耀而行事的人。因为只有那不看重荣耀的人才能获得尊敬和荣耀；但那非常重视大众荣耀，并为获得它而做一切、忍受一切的人，正是那无法获得它的人，并且会遭受与荣耀完全相反的一切：嘲笑、指责、讥讽、敌意和仇恨。这不仅发生在男人中，也发生在你们女人中，而且实际上更特别在你们的案例中。因为那容貌、步态、衣着都不矫饰，不求任何人荣耀的女人，被所有女人钦佩，她们欣喜若狂地赞美她，称她有福，并为她祈求各种各样美好的事物；但一个虚荣的女人，她们以厌恶和憎恨看待她，像躲避某种野兽一样躲避她，并给她加上无尽的咒骂和辱骂。通过拒绝接受大众荣耀，我们不仅躲避了这些邪恶，而且除了已经提到的那些，我们还将获得最高的益处，逐渐被训练放松对土地的把握，向天堂移动，并鄙视所有世俗的事物。因为那不需要来自人的荣耀的人，将安全地行他所行的任何善事，并且在这生命的患难和顺境中都不会受到严重影响；因为前者不能使他沮丧、打倒他，后者也不能使他得意、自高自大，相反，在不确定和动荡的环境中，他本人保持不受任何变化的影响。我预计这很快就会发生在你自己的灵魂上，一旦你永远摆脱所有世俗的利益，你将在我们当中显示一种天堂般的生活方式，不久你就会嘲笑你现在所悲叹的荣耀，鄙视它空洞而虚伪的面具。但是，如果你渴望你以前因丈夫而享有的保障，以及财产的保护，并免受那些践踏他人不幸之人的算计，\BibleQuote{将你的挂虑托付于主，祂必供养你。} \BibleQuote{因为请看，}经上说，\BibleQuote{回顾历代，看谁曾信赖主而蒙羞？或谁曾呼求祂而被忽略？或谁曾持守祂的诫命而被遗弃？}\BibleRef{德 2:10} 因为那已经减轻了这无法承受的灾难，并已使你处于平静状态的那一位，也将避免即将来临的邪恶；因为你自己也会承认，你永远不会再受到比这更沉重的打击。既然你这样勇敢地承受了目前的痛苦，而且在你没有经验的时候，你将更容易忍受将来的事件，如果发生任何违背我们愿望的事（愿天主禁止）。因此，寻求天堂，以及一切有助于来世生命的事物，这里没有任何事物能伤害你，甚至黑暗世界的统治者本人也不能，只要我们不自伤。因为如果有人夺走我们的财产，或将我们的身体砍成碎片，这些事与我们无关，只要我们的灵魂保持完整。\par

7. 现在，我一次性告诉你，如果你希望你的财产安全地与你同在并进一步增加，我将向你展示这个计划，以及那个没有任何图谋它的人能进入的地方。那么，这地方是哪里？是天堂。把你的财产寄给你那善良的丈夫，便没有盗贼、阴谋者或其他任何破坏性的东西能扑向它们。如果你把这些财物存入来世，你会发现从中获得许多利润。因为所有我们在天堂种植的东西都产出大量丰盛的收成，正如那些根在天堂的东西自然可以预料的那样。如果你这样做，看你会享受何等的福分：首先是永生和那些爱天主的人所应许的事物，\BibleQuote{那是眼睛没有看见，耳朵没有听见，人心也未曾想到的}；其次是与你的善良丈夫永久的交往；并且你将摆脱困扰你的忧虑、恐惧、危险、算计、敌意和仇恨。因为只要你被这些财产包围，就很可能有一些人图谋它；但如果你把它转移到天堂，你将过上安全、保障和非常宁静的生活，享受独立与虔诚的结合。因为当一个人想买地，寻找肥沃土地时，如果天堂被提供给他代替地上，并且有在那里获得产业的可能性，他仍然停留在地上，忍受与之相关的辛劳，这是非常不合理的；因为土地常常使我们的希望落空。\par

但既然你的灵魂因常常存着你丈夫会升任总督之位的期望，以及他被过早地从那高位上夺走这一念头而深感忧伤和烦恼，那么首先请考虑这一事实：即使这希望是很有根据的，它也不过是人的希望，常常会落空。我们在生活中看到许多这样的事：那些被满怀信心期待的事未能实现，而那些从未进入脑海的事却常常发生。在政府、王国、继承和婚姻方面，我们随处可见这种情况。因此，即使机会近在眼前，正如俗语所说：\Quote{从杯到唇，常有失手}，而圣经也说：\BibleQuote{从早晨到晚上，时间就改变了。} \BibleRef{德 18:26}\par

同样，今天在这里的国王明天就死了。此外，同一位智者为了说明人希望的逆转，说道：\BibleQuote{许多暴君曾坐在地上，而一个从未被想到的人却戴上了王冠。} \BibleRef{德 11:5} 而且，他若活着，就一定能得到这一高位，这并非绝对确定；因为未来之事是不确定的，使我们产生各种猜疑。凭什么证据表明，他若活着就一定能得到那高位，而不会出现相反的情况，即他会失去他实际拥有的职位，或因患病、或遭那些在顺境中胜过他之人的嫉妒与敌意、或遭受其他某种严重的不幸呢？但让我们姑且假设，他若活着就显然一定能获得这一高位；那么，与这高位的大小相称的是他必须面对的更大危险、焦虑和阴谋。或者将这些都撇开，假设他安全地穿越了那充满困难的海洋，且十分平静；那么告诉我，终点是什么？不是他现在所达到的；不，不是那个，而是别的、很可能是令人不悦和不愿要的东西。首先，他见到天堂和天上的事物将被延迟，这对那些信赖将来之事的人是一个不小的损失；其次，即使他过的是非常纯洁的生活，他生命的长度和他高位上的事务也可能会使他无法像现在这样纯洁地离世。事实上，不确定他是否不会经历许多变化，并在临终前陷入怠惰。因为现在我们确信，靠天主的恩宠，他已飞往安息之地，因为他没有犯过任何排除在天国之外的行为；但在那种情况下，长期接触公共事务后，他可能会沾染很大的污秽。因为在如此大的恶事中行走，要走正直的路是极其罕见的事；而故意或非故意地走歧路则是自然的事，且经常发生。但如今，我们已从这种忧虑中解脱出来，并且坚信在那伟大的日子，他将大放光明，靠近基督闪耀，与天使一起走在基督的前面，穿着无可言喻的荣耀之袍，在基督审判时侍立，并担任基督首要大臣之一。因此，停止哀悼和悲伤吧，你要持守与他同样的生活之道，甚至要更加严谨，使你能迅速达到与他同等的德行标准，居住在同一处所，并永永远远与他重新结合——不是在这婚姻的结合中，而是在另一种远为更好的结合中。因为这只是肉体的交合，但那时将是灵魂与灵魂更完美的结合，一种远为愉悦、远为高尚的结合。\par

\SourceRecord{Translated by W.R.W. Stephens.；Nicene and Post-Nicene Fathers, First Series；Vol. 9.；Edited by Philip Schaff.；Buffalo, NY: Christian Literature Publishing Co.,；1889.；<http://www.newadvent.org/fathers/1904.htm>.}
